\chapter{Conclusion and future work}
\label{chap:conclusion}

\section{Summary of contributions}
\label{sec:conc-summary}

This project tackled the problem of detecting and classifying
radiological anomalies in noisy gamma-dose time series, starting from
four months of real 2015 recordings. The work split into two
iterations:

\paragraph{Iteration 1 (v1)} fitted an IID-Gaussian-plus-Ornstein--
Uhlenbeck noise model on the real data, designed six anomaly templates
and a dense ($\sim 7$ events / $10\text{k}$ samples) synthetic
generator, trained three PyTorch architectures (1D-CNN, ResNet 1D,
CNN+Attention) on a $10^{7}$-sample synthetic corpus, and built the
end-to-end event-level evaluation, robustness sweep, and real-data
inference pipelines that the rest of the project re-uses. On the
synthetic benchmark, all three architectures cleared $90\%$
event-level F1 (the ResNet reaching $100\%$); on the real data, all
three failed in different architecture-specific ways.

\paragraph{Iteration 2 (v2)} produced a deeper background re-analysis
that exposed the v1 noise model's blind spot (the residual is an
AR(1) process with $\varphi \approx 0.86$, not white), hand-labelled
$17$ plausible events across the four months to constrain the
empirical event distribution (durations $[32, 1014]$ samples,
amplitudes $[11, 37]$ counts, only three of the six v1 templates
have an empirical counterpart), and re-designed the synthetic
generator to match all of these targets simultaneously. Retrained on
this corpus, the three architectures sit on the \texttt{Normal} class
on the quiet stretches of the real signal and fire only where a
visible event is present.

\section{Architecture recommendation}
\label{sec:conc-archi}

On the v2 synthetic benchmark the 1D-CNN and the CNN+Attention reach
a perfect event-level macro F1 ($1.00$) and the ResNet sits at $0.92$
on the same $15$-event segment. The gap is too small to act as a
decisive ranking on its own. The separation that matters appears on
two other axes:
\begin{itemize}[leftmargin=*, itemsep=0.2em]
    \item \textbf{Robustness} (\cref{chap:robustness}): the residual
    and attention models degrade much more gracefully under
    background-noise growth than the baseline 1D-CNN, which collapses
    to a single-class plateau already at $m_{\text{noise}} \approx
    1.7$. Between the two robust architectures the ResNet pulls
    slightly ahead (noise $50\%$-crossing at $m_{\text{noise}}
    \approx 2.8$ vs $2.6$ for CNN+Attention, deformation
    $50\%$-crossing at $m_{\text{deform}} \approx 6.4$ vs $5.7$).
    \item \textbf{Real-data behaviour} (\cref{chap:v2real}): the three
    models look qualitatively similar on the v2 real-data plots, but
    the ResNet is the cleanest --- its background confidence stays
    closest to $100\%$ on the quiet stretches without any of the
    residual jitter the 1D-CNN still shows.
\end{itemize}

If a single architecture must be picked for the next iteration, the
recommendation is \textbf{ResNet 1D}: highest robustness, cleanest
real-data inference, and a parameter count ($\sim 75\text{k}$) small
enough to retrain in under ten minutes on a laptop GPU. The single
synthetic misclassification it carries is a $61$-sample asymmetric
\texttt{bell} predicted as \texttt{fast\_ascend} --- not a structural
failure. The CNN+Attention is the second pick: it ties the 1D-CNN at
a perfect synthetic macro F1, comes close to the ResNet on robustness,
and is the most promising candidate for longer-window analyses
(\cref{sec:future-attention}). The baseline 1D-CNN is the fastest to
train but its noise fragility makes it unfit for a beacon that must
run across seasons.

\section{Limitations}
\label{sec:conc-limits}

Several known limitations of the current pipeline are worth recording
explicitly.

\begin{itemize}[leftmargin=*, itemsep=0.25em]
    \item \textbf{No real ground truth.} The four months of real data
    have no expert labelling, so we have no quantitative measure of
    how the v2 models do on real anomalies. The qualitative inspection
    in \cref{chap:v2real} is the best we can do for now.
    \item \textbf{Very short events.} The event-coverage labelling
    rule (\cref{sec:framing}) keeps events of $\sim 30$ samples and
    above inside the training distribution, and the synthetic-test
    segment of \cref{tab:v2-event-prf} detects them on all three
    architectures. Events shorter than this --- $\sim 10$--$20$
    samples, well below the lower bound of the v2 duration range and
    of the hand-labelled real events --- would still drop out of any
    coverage-based labelling at the present stride and window length.
    A finer-grained labelling stage (sample-level or with a much
    smaller window) would be required to address such events, and
    the present pipeline does not.
    \item \textbf{Single beacon.} The pipeline assumes a univariate
    time series. Real surveillance networks involve several beacons
    and several covariates (temperature, humidity, pressure).
    Exploiting these as additional input channels is left for future
    work.
    \item \textbf{Black-box decisions.} The trained models output a
    class label but not a justification. The architectures expose the
    relevant hooks (\texttt{last\_conv\_features}) but a Grad-CAM
    analysis is left for the next iteration.
\end{itemize}

\section{Future work}
\label{sec:conc-future}

\subsection{Expert labelling of real events}

The single highest-value item is expert labelling of the real
recordings. The current $17$ provisional labels are sufficient to
constrain the marginal distributions used by the v2 generator and to
validate the v2 inference qualitatively, but a real test set
--- ideally on more than the four available months, since seasonal
effects on the background are visible in \cref{tab:reanal-drift} ---
would let us replace ``v2 looks better than v1'' by a measured number.

\subsection{Self-supervised pre-training on the real signal}

A natural way to bridge the residual synthetic--real gap is to
pre-train the convolutional stem on the four (and ideally more) months
of real, unlabelled data, using a self-supervised objective such as
masked-segment reconstruction. The stem then learns features adapted
to the real signal statistics; only the classifier head needs to be
fine-tuned on the synthetic dataset.

\subsection{Longer windows and stronger attention}
\label{sec:future-attention}

The CNN+Attention model did not benefit from its long-range mixing on
$512$-sample windows because the v2 templates rarely extend past
$\sim 1\,000$ samples. With longer windows ($W = 4\,096$ or more), the
attention model could in principle disambiguate close events, reason
about context (a slow ramp during which a sharp spike occurs), and
detect multi-event episodes. Scaling the attention up requires the
usual care (longer sinusoidal encoding, more heads or layers, possibly
linear-complexity attention variants) but the architecture is already
in place.

\subsection{An ``unknown but not Normal'' class}

The v2 generator is still binary in spirit: a window either contains a
clean templated event or pure stationary background. Real recordings
contain a richer family of intermediate signals (a slow drift, an
outlier sample, a brief sensor wobble). Introducing a synthetic
template family for ``noise-like but not Normal'' signals (single-sample
spikes, slow ramps below the v2 amplitude floor, sensor dropouts) would
give the network somewhere to put its confidence when the input is not
\texttt{Normal} but does not match any radiological template either.

\subsection{Explainability with Grad-CAM and attention visualisation}

The two convolutional architectures expose their last convolutional
layer as \texttt{last\_conv\_features}, ready to be hooked by a
Grad-CAM implementation. For the attention model, the attention
matrices themselves are interpretable saliency maps over the token
sequence. Adding both to the inference scripts would let the
supervisor visually verify that the model is paying attention to the
right part of an event when it predicts a class.

\subsection{Multi-beacon and multi-modal extensions}

A practical surveillance system never works with a single signal in
isolation; pressure, temperature, humidity and the gamma counts of
neighbouring beacons are all informative. Extending the input to a
small number of channels is mechanical (the convolutions are already
multi-channel by design); the harder problem is collecting a labelled
multi-modal dataset, which the synthetic generator could be extended
to produce.

\section{Final words}
\label{sec:conc-final}

The project ended in a much better place than it started, but the
trajectory is the most informative output. On the benchmark we
designed for ourselves the v1 system worked very well --- all three
architectures cross the $90\%$ event-level F1 line, the residual
network is flawless. On the real data, v1 failed dramatically and in
three architecture-specific ways. The honest negative result, and the
diagnostic infrastructure built around it, are what allowed the v2
re-analysis to be sharp enough to fix the right problem. The single
most useful methodological lesson is therefore: \emph{do not benchmark
a deep model on the same kind of data it was trained on, and call it
done}. The first contact with real data is what revealed where the
pipeline really stood, and the diagnostic infrastructure built to make
that contact diagnosable is, in the long run, the most valuable output
of the project.
